% Introducción del informe final. Las cifras provienen de las secciones de metodología y de resultados.
\section{Introducción}
\label{sec:introduccion}

Retener a un cliente cuesta menos que conseguir uno nuevo, y por eso la predicción de la fuga (\emph{churn}) es
una de las aplicaciones más extendidas del aprendizaje supervisado en los servicios financieros. En una
\emph{fintech} casi toda la relación con el cliente deja huella digital (uso de la aplicación, tenencia de
productos, satisfacción declarada y actividad transaccional), y un modelo capaz de ordenar a los clientes por su
riesgo de abandono permite concentrar una campaña de retención de capacidad limitada en quienes más la
necesitan. Este informe reúne el trabajo de las tres entregas del proyecto final del curso sobre el conjunto
público COFINFAD \citep{cofinfad2025}, que describe a 48.723 clientes de una \emph{fintech} colombiana con 54
columnas y 3,16 millones de transacciones de 2023, y cuyo objetivo, \cp, es un puntaje de propensión al
abandono.

El proyecto avanzó en tres entregas. La primera auditó la calidad de los datos y exploró sus variables; la
segunda preparó los datos y ajustó los modelos lineales de referencia, una regresión logística con penalización
L1 y un Lasso, que alcanzaron un AUC de 0,6832 al clasificar la mitad de mayor riesgo y un $R^2$ de 0,1511 al
predecir el puntaje; la tercera construyó el pipeline combinatorio que pide la guía del curso (siete
clasificadores por cuatro técnicas de balanceo por cuatro métodos de optimización de hiperparámetros, más siete
regresores por cuatro métodos: 140 corridas con validación cruzada anidada de $5 \times 3$ pliegues) y propuso
un modelo nuevo, al que dedicó 20 corridas más. Dos hallazgos de las entregas anteriores condicionan todo el
trabajo. El primero es que el objetivo se generó en buena parte a partir del número de productos activos, \ap,
que por eso se excluye como predictor por fuga de información \citep{kaufman2012leakage}; sin ella existe un
techo de información que se estima y con el que se compara cada modelo. El segundo es que, sin esa variable, la
señal se concentra en cuatro variables (frecuencia de uso de la aplicación, edad y dos medidas de satisfacción)
cuyos efectos actúan por tramos, mesetas y escalones que una recta no puede representar; esa no linealidad
explica el desempeño modesto de los modelos lineales.

Los modelos del curso que sí captan la no linealidad, los ensambles de árboles como Random Forest
\citep{breiman2001rf} y XGBoost \citep{chen2016xgboost}, alcanzan la mejor precisión, pero necesitan
herramientas auxiliares como SHAP \citep{lundberg2017shap} para explicar por qué un cliente queda marcado, y en
un contexto financiero la explicación forma parte del producto. Para cerrar esa brecha se propone una
\emph{Explainable Boosting Machine} (EBM) \citep{lou2012intelligible,lou2013ga2m,nori2019interpretml}, un
modelo aditivo generalizado con interacciones por pares cuyas funciones de forma se aprenden por \emph{boosting}
cíclico de árboles de pocas hojas: la predicción es la suma de curvas que se pueden graficar y auditar. No forma
parte de los modelos del curso. Para adaptarla al problema se implementó la ponderación de clases mediante pesos
por observación y una búsqueda de hiperparámetros multi-fidelidad, y se integró en el mismo pipeline, con las
mismas reglas y el mismo presupuesto, que los modelos de referencia.

La pregunta que responde este informe es si un modelo aditivo legible alcanza la precisión de los ensambles de
árboles en estos datos, y a qué costo. La hipótesis de partida es que sí, porque la señal está en pocas
variables con efectos no lineales y casi independientes, que es exactamente lo que un modelo aditivo
representa. Las contribuciones son cuatro: (i) la EBM adaptada al problema, con su formulación, su pseudocódigo
y la justificación de cada componente; (ii) una comparación con los siete modelos de referencia de cada tarea
bajo un protocolo estadístico fijado de antemano, con pruebas pareadas, corrección por comparaciones
múltiples, regiones de equivalencia práctica y tamaños de efecto; (iii) un estudio de ablación con cinco
hipótesis, análisis de sensibilidad y de costo computacional; y (iv) la interpretación del modelo y el análisis
de sus errores, que localiza dónde falla y por qué. El resultado, anticipado, es que la EBM supera con
significación a los modelos lineales, a k-NN y a Naive Bayes, obtiene el mayor AUC de prueba (0,8289 con la
etiqueta del cuarto superior; en la partición y la etiqueta de la entrega anterior sube el AUC de 0,6832 a
0,7187) y el segundo RMSE, y empata estadísticamente con XGBoost, Random Forest y el árbol de decisión,
cubriendo el 98~\% del techo de información; su inferencia es tan rápida como la de XGBoost, pero su ajuste es
de 15 a 17 veces más lento. La conclusión que la evidencia sostiene es que la EBM es una alternativa equivalente
en precisión y legible por construcción, no que sea mejor.

El informe se organiza así. La sección~\ref{sec:metodologia} formula el problema y describe los datos y su
preparación, el modelo propuesto, su entrenamiento, los modelos de referencia, el protocolo de evaluación y el
diseño de la ablación. La sección~\ref{sec:resultados} presenta el desempeño global, las pruebas de
significancia, la ablación, la sensibilidad, el costo, la interpretabilidad y el análisis de errores, y discute
las limitaciones. La sección~\ref{sec:conclusiones} cierra con la respuesta a las hipótesis y el trabajo
futuro. Toda cifra del informe la producen guiones que se ejecutan sobre los resultados guardados; el código,
los datos y los documentos están en \url{https://github.com/RubenEsg/Prediccion-de-fuga-de-clientes-en-fintech}.
